\section{The \sys{} Method}
\label{sec:method}
% Follows the system figure. Purpose, definition, algorithm, guarantee, complexity for each step.
\sys{} consists of two steps. Step~1 <...>. Step~2 <...>, and its output feeds back into Step~1 for the next round.

\subsection{Step 1: <Name>}
\label{sec:step1}
<Purpose in one sentence>. <Definition>. <How it works on the running example>.

\subsection{Step 2: <Name>}
\label{sec:step2}
<Purpose in one sentence>. Algorithm~\ref{alg:main} lists the procedure.

\begin{algorithm}[t]
\caption{<Name of the procedure>}
\label{alg:main}
\begin{algorithmic}[1]
\REQUIRE candidate set $\cand$, budget $\budget$
\ENSURE labeled set $L$
\STATE $L \gets \emptyset$
\WHILE{$|L| < \budget$}
  \STATE <select the next candidate by the rule of Step~2>
\ENDWHILE
\RETURN $L$
\end{algorithmic}
\end{algorithm}

\begin{theorem}[Guarantee]
\label{thm:guarantee}
Under Assumption <k>, the policy of Algorithm~\ref{alg:main} is no worse than <every baseline it contains>.
\end{theorem}
% State the theorem as a conditional guarantee when the implementation does not meet every premise.
